\documentclass[prep,both]{debate}
\motion{AI的迅猛发展提升 / 降低了人类创作者存在的意义}
\meta{赛制}{招新赛 3v3}
\meta{生成日期}{2026-09-30}
\begin{document}
\debatetitle

\begin{summary}
\begin{fields}
\field{持方优劣评估}{正方较劣势。劣势主要在论据可得性与评委直觉：反方手里是已经发生的订单、收入和丢活的数据，普通人的第一反应也是“AI 在抢创作者的饭碗”；正方的证据多是评价实验、调查意愿和法律规则，行为证据薄。正方的打法是判准之争：意义是需要和珍视两半，反方只量了按件计价的一半；再用“机器越多，人的来源越要被认出来”把价值层做厚，并利用最后结辩的顺序优势。反方要防的，是滑成“就业题”和“AI 取代一切”。}
\thesis{正方}{机器能做的越多，“这是一个人做的”就越被区分、被看重，作品的分量也越落在人的立意上。}
\thesis{反方}{AI 让社会对多数创作者的需要绕开了人，留下来的人手里只剩修补：用 AI，作品里的你变少；不用 AI，需要你的人变少。}
\end{fields}
\end{summary}

\section{一、赛制要点}
\begin{flowtable}
\block{A}
\flow{1 正方一辩立论}{180 秒}{正一}{定义、中立判准、两个论点，抛出全场追问}
\flow{2 反方二辩质询正方一辩}{90 秒 双边}{反二 问}{全场第一次质询：为反一立论开头取材}
\flow{3 反方一辩立论}{180 秒}{反一}{开头接住质询所得，主体是定义、判准、两个论点}
\flow{4 正方二辩质询反方一辩}{90 秒 双边}{正二 问}{拿到“意义包括被珍视”和“法院认人的选择”}
\block{B}
\flow{5 正反双方三辩对辩}{各 90 秒}{正三 反三}{第一个战场：需要和珍视哪个算数}
\flow{6 反方二辩申论}{90 秒}{反二}{全场第一篇申论：拆正方立论主干，推进反方论点一}
\flow{7 正方三辩质询反方二辩}{90 秒 双边}{正三 问}{压数据量级，拿“要求标注是在乎来源”}
\flow{8 正方二辩申论}{90 秒}{正二}{回应反二申论，拆反方立论主干，推进正方论点一}
\flow{9 反方三辩质询正方二辩}{90 秒 双边}{反三 问}{把标识拉回“防骗”，拿“认不出就落不到人身上”}
\block{C}
\flow{10 正反双方一辩对辩}{各 90 秒}{正一 反一}{中段收拢：两半都算，看多数}
\flow{11 正方三辩申论}{90 秒}{正三}{处理反方最强点，推进正方论点二}
\flow{12 反方一辩质询正方三辩}{90 秒 双边}{反一 问}{拿“立意在甲方”和“作品越来越像”}
\flow{13 反方三辩申论}{90 秒}{反三}{回应正三申论，处理正方最强点，推进反方论点二}
\flow{14 正方一辩质询反方三辩}{90 秒 双边}{正一 问}{全场最后一次质询，为正一结辩取材}
\block{D}
\flow{15 正反双方二辩对辩}{各 90 秒}{正二 反二}{结辩前最后一次交锋，守住判准}
\flow{16 反方一辩结辩}{90 秒}{反一}{先结辩：归纳分歧、处理最强点、封路}
\flow{17 正方一辩结辩}{90 秒}{正一}{最后一发：回应反一结辩，留约 40 字临场位}
\end{flowtable}
质询全部是双边计时，频繁打断要扣分，被质询方故意不答也扣分：两条链共四到七问，问题短到一句话能答完，对方答非所问先复述一次再问。没有驳论、小结和自由辩，对方立论只能在两篇申论里拆，所以两篇申论赛前分工写死（见口径表）：二辩申论拆对方立论主干并推进我方论点一，三辩申论处理对方最强点并推进我方论点二。一辩稿算全队的作品，一辩的个人分来自被质询、一辩对辩、质询对方三辩和结辩，这几处材料要给足。

\section{二、辩题性质与争议焦点}
\begin{fields}
\claim[辩题性质]{比较型的事实判断加价值判断。“提升 / 降低了”是和 AI 迅猛发展之前比的变化方向，用的是完成时，看已经发生的事和已经显现的趋势。}
\field{正方倾向把题目解释成}{价值题：人的来源、人的立意在社会上的分量是否更重。重点放在“被珍视”和“作品里人的位置”。}
\field{反方倾向解释成}{处境题：多数创作者被需要的程度、作品里属于自己的部分是否变少。重点放在“被需要”和“创作过程”。}
\field{争议焦点}{\sub{需要和珍视哪个算数}{两方都会同意两半都算，真正的交锋是珍视落在谁身上、需要在谁身上缩水。}\sub{修图工还是立意者}{AI 接走执行以后，人是升级成决定画什么的人，还是降级成给机器收尾的人。}\sub{看多数}{正方的证据多是规则与评价，反方的证据多是收入与订单，谁更贴近“多数创作者的常态”。}}
\end{fields}

\section{三、关键词定义}
\subsection{3.1 AI 的迅猛发展}
\begin{fields}
\field{调研}{\sub{常识义}{以 2022 年底的 ChatGPT、同年的 Midjourney 与 Stable Diffusion 为标志，生成式人工智能在文本、图像、音乐、声音上的能力与普及度快速提升（常识，无需单独出处）。}\sub{法规义}{《人工智能生成合成内容标识办法》把“人工智能生成合成内容”作为需要标识的对象，2025 年 9 月 1 日起施行（网信办等四部门）。}}
\begin{procard}{正方有利定义}\definition{AI 的迅猛发展，是创作工具能力的迅速扩张。}它把 AI 定位成给人用的笔。\sub{有利之处}{替代问题被淡化，论点二“立意回到人”顺理成章。}\sub{反方如何反驳它}{生成式 AI 可以不经过人直接产出成品，所以才有 Thaler 案和标识办法；只把它叫“工具”，是回避题目里最尖锐的那一面。}\end{procard}
\begin{concard}{反方有利定义}\definition{AI 的迅猛发展，是机器可以绕过人独立产出作品的能力。}\sub{有利之处}{直接对接“不再被需要”。}\sub{正方如何反驳它}{独立产出的东西在法律上连作品都不算（Thaler 案；《著作权法实施条例》第三条把创作定义为智力活动）；AI 的发展同时包括人机协作，只取一面同样偏。}\end{concard}
\begin{neutralcard}{中立定义}2022 年以来，机器生成文字、图像、音乐和声音的能力突然变强、人人可用；它既能被人拿来用，也能不经过人直接产出。\sub{合理性}{贴近常识，两面都保留，双方都有得打。}\end{neutralcard}
\end{fields}

\subsection{3.2 人类创作者}
\begin{fields}
\field{调研}{\sub{词典义}{汉典（《国语辞典》）“创作”：“特指不事模仿、出于己意的创造。多指文学、艺术等作品的创造而言。”}\sub{法规义}{《著作权法实施条例》第三条：“著作权法所称创作，是指直接产生文学、艺术和科学作品的智力活动。”第二款：组织、咨询、提供物质条件等辅助工作“均不视为创作”。}}
\begin{procard}{正方有利定义}\definition{凡是出于己意创造作品的人都是创作者，包括业余者和借 AI 表达的人。}\sub{有利之处}{把“更多人能创作”划进来，也让“甲方也是创作者”成立。}\sub{反方如何反驳它}{把“创作者”稀释成“所有发帖的人”，躲开了以创作为志业的那群人正在失去的东西；题目说的是“人类创作者”，不是“人类”。}\end{procard}
\begin{concard}{反方有利定义}\definition{以创作为业、靠作品谋生的职业创作者。}\sub{有利之处}{替代与收入数据直接落在他们身上。}\sub{正方如何反驳它}{意义不是工资单，一个没有稿费的诗人也有存在的意义；把业余者排除，等于把题目改成“AI 降低了创作行业的收入”。}\end{concard}
\begin{neutralcard}{中立定义}把创造文艺作品当成自己一部分的人，靠它吃饭的算，业余在写在画的也算。\sub{合理性}{贴近“出于己意的创造”这一词典义；不以收入划线，也不把偶尔下单的人算进来；双方各有空间：正方讲业余者的表达，反方讲不出名的大多数。}\end{neutralcard}
\end{fields}

\subsection{3.3 存在的意义（以及“提升 / 降低”）}
\begin{fields}
\field{调研}{\sub{词典义}{汉典“意义”：“事物所包含的思想内容和道理”；另一义项“作用与重要性”（英译 value）。本题取后者。}\sub{比较对象}{“提升 / 降低了”是和同一群人在 AI 迅猛发展之前的状况比。}}
\begin{procard}{正方有利定义}\definition{存在的意义，是人类创作者身上机器替代不了的那部分价值：立意、经历、责任。}\sub{有利之处}{AI 越强，这部分越凸显。}\sub{反方如何反驳它}{这是把“剩下的”定义成“全部的”：缩水之后说纯度更高，是循环论证。照这个定义，AI 替代什么都不算损失。}\end{procard}
\begin{concard}{反方有利定义}\definition{存在的意义，是社会对人类创作者的实际需要：订单、岗位、收入、关注。}\sub{有利之处}{数据全在这一边。}\sub{正方如何反驳它}{把意义等同于市场价格；照这个定义，梵高生前没有存在的意义，这不是常识。}\end{concard}
\begin{neutralcard}{中立定义}人类创作者的作用与重要性，分两层：对外，这个世界还有多需要、多珍视由人来创作；对内，创作还能给创作者自己多少东西。\sub{合理性}{直接取自词典义；需要和珍视都算，谁也不能只看一层。}\end{neutralcard}
\end{fields}

\subsection{3.4 隐含要素}
\begin{fields}
\field{主体}{人类创作者整体，重点看多数：不只看顶尖少数，也不只看被淘汰的个人。}
\field{范围}{全球，含中国；覆盖文字、图像、音乐、声音四类。}
\field{时间尺度}{2022 年以来到今天，加上已经显现的趋势；不押注遥远的未来。预测类数据（如到 2028 年）只能说趋势。}
\field{比较对象}{同一群人在 AI 迅猛发展之前的状况，不是和“没有任何技术”的世界比。}
\end{fields}

\begin{warnbox}{3.5 定义杀陷阱：双方都要背}
\begin{fields}
\begin{procard}{正方的定义杀陷阱}把意义定义成“只有人能给的东西”，于是 AI 替代什么都不算损失。评委会认为正方在回避替代的事实。正方必须主动承认执行型订单在减少，再把战场拉到两半。\end{procard}
\begin{concard}{反方的定义杀陷阱}把意义定义成“收入与岗位”，于是只要有人失业就赢。评委会认为反方把意义降成了工资单。反方必须承认珍视那一半存在，再论证它落在少数人身上。\end{concard}
\end{fields}
\end{warnbox}

\section{四、判准与论证义务}
\begin{fields}
\claim[判准是什么]{和 AI 迅猛发展之前比，由人来创作这件事的作用与重要性，整体上是变重还是变轻。}
\begin{procard}{正方有利判准}\definition{看“人之所以为创作者”的那部分（立意、经历、判断）是否更被看见、更被珍视。}\sub{有利之处}{AI 越强，这部分越凸显。}\sub{反方如何反驳它}{只看剩下的最好的部分，回避被替代掉的大多数；“更被珍视”的可能只是一小撮有名字的人，这是拿幸存者证明群体。}\end{procard}
\begin{concard}{反方有利判准}\definition{看社会还需要多少人类创作者：订单、岗位、收入是增是减。}\sub{有利之处}{替代数据直接成立。}\sub{正方如何反驳它}{需要不等于意义；按件计价的需要下降，不等于人被珍视的程度下降。照这个标准，摄影术出现后画家的意义就降低了。}\end{concard}
\begin{neutralcard}{中立判准}看两层，而且看多数：对外，社会对“由人来创作”的需要与珍视，整体更重还是更轻；对内，创作者从创作里得到的东西，更多还是更少。说的是创作者群体的常态，不是顶尖几个人，也不是个别被淘汰者。\sub{合理性}{贴近常识：意义既是被需要，也是被珍视；看多数，防止双方各举极端例子。}\sub{正方需要证明}{AI 之后，人来创作这件事在社会上被更需要或更珍视，而且这种变化不只发生在顶尖少数；创作者本人从创作里得到的东西没有整体变少，或有具体机制让它变多。}\sub{反方需要证明}{AI 之后，人来创作的必要性与被珍视程度加在一起整体变轻，被替代与被掏空的是多数创作者的常态，而不只是边缘的一部分。}\end{neutralcard}
\field{义务提醒}{两方都要证明整体方向，都不能只证明“有”：正方不能只证“有人更被珍视”，反方不能只证“有人丢了活”。一辩稿都用中立判准。}
\end{fields}

\prosection{五、正方论点}
\prosubsection{论点一：人味成了稀缺品}
\begin{fields}
\claim{机器内容无限以后，社会开始专门区分人和机器，并且在能选的时候选人。}
\begin{mechanism}\step AI 把作品的数量推向无限，成本推向零。\step 稀缺的不再是作品，而是作品背后那个有经历、要负责的人。\step 社会用评价、标注和法律来区分并看重人的来源。\end{mechanism}
\begin{evidence}{学理}[已核实]原作价值的“表演说”（Newman \& Bloom，2012）\sub{核心主张}{人们看重原作，是因为把作品当成一次独特的创作行为来估值，而不只看物件本身。}\sub{支撑}{机制第二步：人们估值的对象是“谁做的、怎么做的”。}\sub{边界}{它说明人怎么估值，不保证市场一定付钱。旁证“努力启发式”（Kruger 等，2004）的复制结果混合，出价那一项没有复制出来，只作边界，不作主证据。}\sub{口头版}{我们看重一件作品，是看重一次创作行为。}\sub{出处}{Newman, G. E. \& Bloom, P. (2012). Art and authenticity: the importance of originals in judgments of value. JEP: General. DOI 10.1037/a0026035。}\end{evidence}
\begin{evidence}{数据}[已核实]同一批由 AI 画的画，随机贴上“人类创作”或“AI 创作”标签；贴人类标签的，喜欢、美感、深刻、价值四项评分全部更高。\sub{来源}{Bellaiche, L. 等 (2023). Humans versus AI: whether and why we prefer human-created compared to AI-created artwork. Cognitive Research: Principles and Implications. DOI 10.1186/s41235-023-00499-6。}\sub{方法}{实验操纵标签、不操纵画作，两项研究；第二项加测叙事感与努力感，发现二者调节标签效应。}\sub{局限}{测的是评价，不是购买；画作全部由 AI 生成，反方会说人们看重的是一张标签。}\end{evidence}
\begin{evidence}{数据}[已核实]皮尤研究中心 2025 年调查 5,023 名美国成年人：76\% 认为能分辨图片、视频和文字出自 AI 还是人“极其或非常重要”。Deezer 与益普索 2025 年在 8 国调查 9,000 人：80\% 认为纯 AI 音乐应清楚标注，52\% 因分不出人和 AI 感到不安。\sub{来源}{Pew Research Center, 2025-09-17, How Americans View AI and Its Impact on People and Society；Deezer Newsroom, 2025-11-12, Deezer/Ipsos 调查。}\sub{方法}{两项都是问卷调查；前者为美国成人样本，后者为 8 国样本并含盲听测试。}\sub{局限}{测的是意愿，反方会说“想分辨是怕被骗”。同一份益普索调查里 97\% 的人盲听分不出，是反方的数据。}\end{evidence}
\begin{evidence}{案例}[已核实]制度正在把“作者是人”写进规则：美国哥伦比亚特区巡回上诉法院 2025 年 3 月 18 日在 Thaler v. Perlmutter 案判定版权法上的作者必须是人，最高法院 2026 年 3 月 2 日拒绝调卷；我国《人工智能生成合成内容标识办法》2025 年 9 月 1 日起施行。\sub{代表性}{是规则而不是个例，对所有创作者一样适用，回应“看多数”。}\sub{出处}{Authors Alliance 2025-03-19 对判词的引述；National Constitution Center 2026；网信办 2025-03-14。}\end{evidence}
\closing{对外层的“珍视”变重了，而且标注和法律不挑名气，对每一个人类创作者都适用。}
\begin{clash}\pair{被区分不等于被需要，标签保护的是不被冒名，不是有人来请。}{我方说的是珍视，不是价格；被认出来是被选择的前提，认不出的东西谁也选不了。}\pair{保护恰恰说明濒危。}{意义是重要性，不是安全感；人只会为自己看重的东西立法。}\pair{标识办法写明是为了防虚假信息。}{目的是什么都好，结果是“是不是人做的”成了必须公开的事实。正方不说它“为了尊重创作者”。}\end{clash}
\end{fields}

\prosubsection{论点二：立意回到人}
\begin{fields}
\claim{AI 接走“画得像、写得顺”，作品的分量落在人的立意与选择上；门槛降下来，更多有话要说的人成了创作者。}
\begin{mechanism}\step 执行成本被机器压低。\step 作品之间的差别越来越取决于想做什么、选哪一张。\step 创作者的价值集中到立意上，更多人能凭立意进入创作。\end{mechanism}
\begin{evidence}{学理}[已核实]“生成式联觉”（Zhou \& Lee，2024）\sub{核心主张}{人负责探索和筛选，机器负责执行；立意与筛选是用 AI 创作的必要技能。}\sub{支撑}{机制第二、三步。}\sub{边界}{样本是主动使用 AI 的绘画平台用户，不能直接推到所有职业画师。}\sub{口头版}{人出主意，机器动手，分数给了主意。}\sub{出处}{Zhou, E. \& Lee, D. (2024). Generative artificial intelligence, human creativity, and art. PNAS Nexus. DOI 10.1093/pnasnexus/pgae052。}\end{evidence}
\begin{evidence}{数据}[已核实]追踪 5 万多名用户的 400 万件作品：用上文生图 AI 后产量提高 25\%，每次浏览获赞的概率提高 50\%；能成功探索新想法的人被同行评价更高；获赞在采用者之间的集中度下降。\sub{来源}{Zhou \& Lee (2024), PNAS Nexus。}\sub{方法}{平台数据，比较用户采用 AI 前后的产出与同行反馈。}\sub{局限}{自选样本；同一研究还发现平均内容新颖度下降、视觉新颖度全面下降，这是反方的数据，正方必须主动承认“平均在趋同”。}\end{evidence}
\begin{evidence}{数据}[已核实]在线写作实验：获得 AI 故事点子的人，故事被评为更有创意、写得更好，原本最不会写的人进步最大；但用 AI 的故事彼此更相似。\sub{来源}{Doshi, A. R. \& Hauser, O. P. (2024). Generative AI enhances individual creativity but reduces the collective diversity of novel content. Science Advances. DOI 10.1126/sciadv.adn5290。}\sub{方法}{随机对照的在线实验，第三方评价故事。}\sub{局限}{短故事、实验情境；趋同一条是反方的。}\end{evidence}
\begin{evidence}{案例}[已核实]北京互联网法院“AI 文生图”著作权案（2023，李某诉刘某）：法院认定原告设计提示词、设置参数，对画面元素和布局作了个性化的选择和安排，“能够体现出原告的独创性智力投入”，享有著作权。\sub{代表性}{是裁判规则，对所有使用者适用。}\sub{出处}{北京市人大常委会网站 2024-01-23 转北京高院工作报告解读。}\end{evidence}
\begin{evidence}{案例}[已核实]美国乡村歌手 Randy Travis 2013 年中风后几乎失语；2024 年 5 月，团队用 AI 从他的 42 段旧录音中还原嗓音，由他本人和老制作人把关，发行新歌《Where That Came From》。\sub{代表性}{个例，只用来让机制被看见，不证明普遍。}\sub{出处}{CBS News，2024-05-05。}\end{evidence}
\closing{对内层，创作者的立意被实现得更充分；对外层，作品的分量落在人的选择上；获赞集中度下降说明好处不只给头部。}
\begin{clash}\pair{立意回到的是甲方，画师只剩修图。}{出于己意去创造的人就是创作者，立意在谁手里都在人手里；被压价的是执行，想画什么机器替不了。}\pair{平均新颖度在降，作品越来越像。}{越趋同，有立意的人越显眼；同一研究里分数给了探索新想法的人。}\pair{Randy Travis 是明星个例。}{是个例，只说明机制；普遍性交给 5 万多名用户的数据和法院规则。}\end{clash}
\end{fields}

\subsection{（被淘汰的正方候选论点，交给反驳库备用）}
\begin{enumerate}[leftmargin=1.8em]
\item 镜子效应：AI 逼人追问什么是人的创作。没有数据支撑，淘汰。
\item 内在快乐不变：创作的乐趣不因 AI 改变。只能证明“不降”，证明不了“提升”，淘汰。
\item 摄影术类比：摄影没有杀死绘画。“这次不一样”一句就能打掉，降为备用案例。
\item AI 催生新岗位：提示词工程、AI 美术指导。数据弱，容易被“新岗位招的不是原来那群人”打掉。
\end{enumerate}

\consection{六、反方论点}
\consubsection{论点一：不再被需要}
\begin{fields}
\claim{对大多数不出名的创作者，被需要就是他存在意义里最实的一块；AI 让这份需要大面积绕开了人，连最好的那批也挡不住。}
\begin{mechanism}\step 多数创作者靠功能性创作被需要：插画、文案、翻译、配音。\step AI 以接近零的成本满足同样的需要。\step 订单和收入下降，人被需要的理由变少，而且质量换不来保护。\end{mechanism}
\begin{evidence}{学理}[已核实]超级明星经济（Rosen，1981）\sub{核心主张}{当技术让少数人能同时服务所有人，收入会集中到顶端少数。}\sub{支撑}{解释为什么“人的溢价”即使存在，也只落在能被认出名字的少数人身上。}\sub{边界}{它讲的是分配，不直接证明总量下降；正方会拿“获赞集中度下降”反驳，反方要说那是用 AI 的网站用户内部。}\sub{口头版}{溢价只流向认得出名字的少数人。}\sub{出处}{Rosen, S. (1981). The Economics of Superstars. American Economic Review 71(5): 845--858。}\end{evidence}
\begin{evidence}{数据}[已核实]ChatGPT 发布后，一个大型在线自由职业平台上，写作类自由职业者月订单下降 2\%、月收入下降 5.2\%；图像类在图像生成模型发布后月订单下降 3.7\%、收入下降 9.4\%；过去表现越好的顶尖自由职业者受冲击反而更大。\sub{来源}{Hui, X., Reshef, O. \& Zhou, L. (2024). The Short-Term Effects of Generative Artificial Intelligence on Employment: Evidence from an Online Labor Market. Organization Science 35(6). DOI 10.1287/orsc.2023.18441。}\sub{方法}{双重差分：比较受影响职业与不受影响职业在模型发布前后的变化。}\sub{局限}{短期、单一平台、只看自由职业者；正方会压量级（2\%）。}\end{evidence}
\begin{evidence}{数据}[已核实]英国作家协会 2024 年 1 月调查，787 份回应：26\% 的插画师、36\% 的译者已经因为生成式 AI 丢了活；65\% 的小说作者认为 AI 会损害未来收入。CISAC 委托 PMP Strategy 的研究预测：到 2028 年，音乐与视听创作者分别有 24\% 和 21\% 的收入面临风险，配音与字幕译者为 56\%。\sub{来源}{Society of Authors 2024 调查，经欧洲作家理事会（EWC）全文转载核实；CISAC 2024-12 新闻稿。}\sub{方法}{前者为会员自愿问卷；后者为案例研究、市场渗透估计与行业访谈得出的预测。}\sub{局限}{前者自报、自选样本；后者是预测，且 CISAC 代表创作者利益，必须说成“预测”。}\end{evidence}
\begin{evidence}{案例}[已核实]全国首例 AI 生成声音人格权侵权案：配音师殷某的声音被 AI 化后做成文本转语音产品出售，北京互联网法院 2024 年 4 月 23 日一审判赔 25 万元。杭州一位游戏行业猎头对 Rest of World 说，插画岗位一年内降了约 70\%。\sub{代表性}{前者是判决事实，说明替代已经发生到“声音”这一层；后者是单个猎头的估计，只能作画面，不能说成行业统计。}\sub{出处}{中新网 2024-04-25；Rest of World，Viola Zhou，2023-04-11。}\end{evidence}
\closing{对外层的“需要”变轻了，而且落在多数不出名的创作者身上；顶尖者受冲击更大，说明质量换不来被需要。}
\begin{clash}\pair{意义不等于订单，反方只算了一半。}{两半都算，我方承认；但珍视落在少数有名字的人身上，需要在多数人身上缩水，净额是变轻。}\pair{月订单只降 2\%。}{收入降得更多，图像类将近一成；而且过去评价越好的人受冲击越大。}\pair{CISAC 是利益相关方的预测。}{是预测，我方只用它说趋势；已经发生的看平台研究与作者调查，方向一致。}\end{clash}
\end{fields}

\consubsection{论点二：手艺被掏空}
\begin{fields}
\claim{留下来的人，作品里属于他自己的那部分正在变少：构思去了甲方，执行给了机器，创作者手里只剩修补。}
\begin{mechanism}\step AI 让甲方可以直接生成初稿。\step 创作者被请来修光线、修手指。\step 作品里属于他的风格与判断越来越少，作品之间越来越像。\end{mechanism}
\begin{evidence}{学理}[有把握]劳动过程理论（Braverman，1974）\sub{核心主张}{生产被拆成“构思”和“执行”两半，构思集中到管理方，执行者被去技能化。}\sub{支撑}{机制第一、二步：AI 把构思交给甲方、把执行交给机器。}\sub{边界}{去技能化命题在学界有争议；本次只核到出版信息，没有核到原句。}\sub{口头版}{构思去了甲方，执行给了机器。}\sub{出处}{Braverman, H. (1974). Labor and Monopoly Capital. Monthly Review Press。}\end{evidence}
\begin{evidence}{数据}[已核实]同一项 5 万多名用户、400 万件作品的研究：采用 AI 后，平均内容新颖度下降，视觉新颖度（像素级的风格）峰值和平均值都持续下降。写作实验：用了 AI 点子的故事彼此更相似。\sub{来源}{Zhou \& Lee (2024), PNAS Nexus；Doshi \& Hauser (2024), Science Advances。}\sub{方法}{见正方论点二。}\sub{局限}{样本是 AI 用户和实验参与者；峰值内容新颖度在上升，是正方可用的一面。}\end{evidence}
\begin{evidence}{案例}[已核实]插画师 Amber Yu 曾花整整一个星期画一个穿传统服装舞狮的女子，一张游戏海报收 3,000 到 7,000 元；游戏公司改用 AI 后，只请她修光线和画歪的肢体，价钱是原来的十分之一。重庆一家工作室 15 名角色设计师裁掉 5 名，“两个人能干原来十个人的活”；广东一位游戏画师说 AI 让她们“更高效，也更累”。\sub{代表性}{个人经历，用来让机制被看见；普遍性交给新颖度数据。}\sub{出处}{Rest of World，Viola Zhou，2023-04-11，AI is already taking video game illustrators' jobs in China。}\end{evidence}
\closing{对内层，创作者从创作里得到的东西变少了。两点合成夹击：用 AI，作品里的你变少；不用 AI，需要你的人变少。}
\begin{clash}\pair{创作的本义是出于己意，执行从来不是意义本身。}{风格、笔触、句子的节奏，就是“己意”在作品里的样子；视觉新颖度全面下降，说的正是这部分在消失。}\pair{甲方也是创作者，立意还在人手里。}{立意在人手里，不在创作者手里；题目问的是人类创作者。}\pair{只有用 AI 的人才会趋同。}{不用 AI 的人订单在少，那是论点一；两点合起来，两头都在变轻。}\end{clash}
\end{fields}

\subsection{（被淘汰的反方候选论点，交给反驳库备用）}
\begin{enumerate}[leftmargin=1.8em]
\item 看不见的作者：作品被 AI 洪流淹没。被 Deezer“上传过半、播放只占 1--3\%”正面打穿，降为申论与对辩素材。
\item 认不出是人：普通人分不出人和机器。标识制度正在补，单独立论会被“执行问题”打掉；并入论点一的防守和反驳库。
\item 作品沦为训练原料：谈的是训练公平，不是意义；正方可用赔偿与判决反打，淘汰。
\item 学徒路断：入门订单消失，下一代长不出来。缺直接数据，淘汰。
\item 创作者更焦虑：只能作证据层，淘汰。
\end{enumerate}

\prosection{七、正方反驳库（针对反方全部可能论点）}
\begin{rebuttals}
\rebuttal{1}{不再被需要：订单和收入在下降，连顶尖的人也挡不住。}{订单少了，说明机器能做的活变多了，不说明人做的东西变轻了。}{打判准：意义是需要和珍视两半，反方只算了按件计价的那一半。打论据：最常被引的平台研究里写作类月订单降 2\%；英国作家协会是 787 份自愿问卷；CISAC 是收版税机构的预测。量级和性质都撑不起“整体变轻”。}{“26\% 的插画师丢了活不是事实吗？”→ 是事实，我方承认执行型订单在减少；丢一份活，不等于人来创作这件事变轻。}{最终}
\rebuttal{2}{手艺被掏空：人被压成给机器收尾的修图工。}{机器接走的是手上的重复，留给人的是决定画什么。}{打定义：创作是“出于己意的创造”，核心是己意，不是工时。打机制：插画师被压价是执行降价；北京互联网法院把作品判给做选择的人；Zhou \& Lee 同一研究里，分数给了探索新想法的人。}{“视觉新颖度全面下降呢？”→ 那是平均在趋同；越趋同，有立意的人越显眼。}{最终}
\rebuttal{3}{作品被淹没：每天约九万首 AI 曲目，超过新上传一半。}{上传多不等于被听见。}{打论据：同一平台 AI 曲目只占播放 1--3\%；部分原因是平台主动不推荐 AI 曲目，理由是减少对艺人收入的冲击。这本身就是平台在把人和机器分开。}{“平台干预不是听众的选择”→ 对，正是平台、法律、行业在主动把人和机器分开，这就是珍视。}{淘汰}
\rebuttal{4}{认不出：97\% 的人盲听分不出，46.6\% 的准确率低于瞎猜。}{分不出来的人里一半以上感到不安，八成要求标注。}{打机制：分不出来说明机器学得像，不说明人不重要；Bellaiche 实验显示一旦知道是人做的，评价全面更高；认不出是信息问题，标识办法已在施行。}{“标注了也没人买”→ 我方说的是珍视，不是价格。}{常见}
\rebuttal{5}{同质化：用 AI 的故事越来越像。}{机器让作品越像，有自己立意的人越显眼。}{打论据：同一研究里最不会写的人进步最大；趋同是用同一套点子的风险，解药正是人的立意。}{“集体多样性下降是损失”→ 是风险，我方承认；风险恰恰说明差异只能来自人。}{淘汰}
\rebuttal{6}{作品沦为训练原料，被白拿。}{训练公平是分配问题，不是意义问题。}{打判准：题目比的是意义；而且规则在向人倾斜：WGA 合同保留主张训练违法的权利，Thaler 判词把版权目的落在作者的劳动。}{“被白拿不是贬低吗？”→ 机器离不开人的作品，恰恰说明人替代不了。}{淘汰}
\rebuttal{7}{学徒路断：入门订单没了，下一代长不出来。}{入门的路变了，不是没了。}{打论据：写作实验里新手进步最大；获赞集中度下降，新人更有机会被看见。}{“没订单怎么养学徒”→ 那是职业结构问题，反方要证明的是意义整体变轻。}{淘汰}
\rebuttal{8}{创作者更焦虑：65\% 的小说作者认为 AI 会损害未来收入。}{担心收入，不等于觉得创作没意义。}{打论据：问卷问的是“未来收入”，不是“创作还有没有意义”，把收入焦虑当成意义下降是偷换。}{“焦虑本身就是意义感在下降”→ 焦虑说明在乎；正因为创作重要，才怕失去。}{淘汰}
\rebuttal{9}{AI 画得比人好，人没必要再创作了。}{创作从来不是比谁画得像。}{打定义：出于己意的创造；Newman \& Bloom：作品被当成一次创作行为估值。相机比画家画得像，不是今天才发生的事。}{“那人为什么还要学画？”→ 因为要说的是自己的话，机器替不了“我想说什么”。}{常见}
\end{rebuttals}

\consection{八、反方反驳库（针对正方全部可能论点）}
\begin{rebuttals}
\rebuttal{1}{人味成了稀缺品：社会在区分人和机器，并且选人。}{被区分不等于被需要，标签保护的是不被冒名，不是有人来请。}{打判准：看多数，稀缺品的溢价只流向能被认出名字的少数（Rosen）。打论据：Bellaiche 的画全是 AI 画的，人们看重的是标签；益普索调查里 97\% 盲听分不出，被看重的前提是被认出来。}{“76\% 想分辨，不是在乎人？”→ 想分辨是怕被骗；在乎，和请你画，是两回事。}{最终}
\rebuttal{2}{立意回到人：执行交给机器，价值落在立意上。}{立意回到的是甲方，不是画师。}{打机制：构思与执行分离（Braverman）；Amber Yu 只被请来修图，价钱是原来的十分之一。打论据：Zhou \& Lee 是主动用 AI 的网站用户，且平均内容新颖度下降、视觉新颖度全面下降。}{“甲方也是创作者”→ 那就承认了：把创作当志业的那群人，被下单的人替换了；题目问的是他们。}{最终}
\rebuttal{3}{法律只认人是作者：Thaler 案、北京互联网法院。}{机器不能署名，不等于有人来请你。}{打机制：北京互联网法院认的独创性已经缩到提示词和参数；标识办法答记者问写明是为了提醒用户辨别虚假信息。}{“法律保护就是珍视”→ 被单独保护，说明处境变差；一道防护栏，不是一块奖牌。}{常见}
\rebuttal{4}{门槛降低，更多人能创作。}{人人都能生成，“创作者”这个身份越来越轻。}{打论据：Doshi \& Hauser 同一研究：个人变好、集体趋同；Deezer 每天约九万首 AI 曲目，变多的是作品，不是被需要的作者。}{“表达变容易不好吗？”→ 好，我方不反对表达；但这不等于以创作为志业的人更重要。}{常见}
\rebuttal{5}{镜子效应：AI 逼人重新思考什么是人的创作。}{被迫自证是人，本身就是分量在变轻。}{打论据：Authors Guild 推出人类创作认证，理由正是 AI 生成的书“越来越像人写的”。}{“思考本身有价值”→ 有，但题目问的是创作者，不是哲学家。}{淘汰}
\rebuttal{6}{摄影术类比：摄影没有杀死绘画。}{摄影只会拍，AI 会画、会写、会唱、会配音。}{打论据：替代面不同，CISAC 预测配音字幕译者 56\% 的收入面临风险；Hui 等的研究同时覆盖写作与图像。}{“每次技术革命都有人喊狼来了”→ 这次不是喊，是数据：顶尖者受冲击更大。}{淘汰}
\rebuttal{7}{创作的内在快乐不因 AI 改变。}{不变不等于提升，正方要证的是提升。}{打判准：论证义务在正方；而且内在体验也在变，广东画师说 AI 让她们“更高效，也更累”。}{“你还能自己画呀”→ 能画，但快乐和意义不是一回事。}{淘汰}
\rebuttal{8}{AI 催生新岗位：提示词工程师、AI 美术指导。}{新岗位招的是会用机器的人，不是原来那群画师。}{打论据：重庆工作室 15 名角色设计师裁掉 5 名，“两个人能干原来十个人的活”。}{“转型的人意义不就回来了？”→ 转型成修图，就是我方第二点。}{淘汰}
\rebuttal{9}{AI 只是工具，就像 Photoshop。}{Photoshop 不会自己画完一张海报，AI 会。}{打定义：生成式 AI 可以不经过人直接产出，所以才有 Thaler 案和标识办法。}{“最终还是人在用”→ 用的人从画师变成了甲方，这就是我方第一点。}{常见}
\end{rebuttals}

\section{九、持方优劣评估}
\begin{assessment}
\assess{定义空间}{“意义”在词典里是作用与重要性，常识里既包括被需要也包括被珍视，双方都能在中立定义里找到空间；正方的“稀缺即珍贵”需要多解释一步，反方的“被需要”一听就懂。}{均衡}
\assess{论证义务}{两方都要证明整体方向。正方要证明“提升”，也就是有东西在变重，而且不只发生在顶尖少数；反方只要证明加总变轻。正方还必须先承认执行型订单在减少，再证明另一半的增长盖过它，义务更重。}{略偏反}
\assess{论据可得性}{反方有已经发生的订单、收入、丢活数据（Hui 等 2024、英国作家协会 2024）和生动的中国案例；正方的证据多是评价实验、意愿调查和规则，行为证据薄，Deezer 播放份额又被平台政策污染，不能用。}{偏反}
\assess{评委直觉}{没有评委信息，只作一般判断：普通人的第一反应是“AI 在抢创作者的饭碗”，正方要先克服直觉；但辩论评委通常欣赏“稀缺让人更重要”这类反直觉而有机制的论证。}{略偏反}
\assess{赛制顺序}{正方先立论、最后结辩，有最后一发；反方先申论、先结辩，可以封路。招新赛没有自由辩，最后一发的价值更高。}{略偏正}
\end{assessment}
\begin{fields}
\claim[结论]{正方较劣势，劣势主要来自论据可得性与论证义务：反方的数据是“已经发生”，正方的数据是“人们怎么看”。}
\begin{procard}{劣势方打法}判准之争：一辩稿就把“需要和珍视两半都算”立死，全场追问反方珍视那一半算过没有。收窄战场：主动承认执行型订单在减少，把比赛拉到“人的来源与立意”的分量上，不在订单数字上硬拼。价值层：结辩升华落在“一个想说话的人身后还站着一个我”。顺序利用：正一结辩是全场最后一发，回应组按反一可能的收法写好，再留临场位。\end{procard}
\begin{concard}{优势方要防的}滑成“就业题”：只讲订单与收入，被正方一句“意义不是工资单”打成偷换。说过头：把“降低”说成“AI 取代了创作者”，被正方用法律和标注反证。把 CISAC 的预测说成已经发生、把猎头的“约七成”说成行业统计，被质询一问就倒。\end{concard}
\end{fields}

\section{十、口径表}
\prosubsection{10.1 正方口径表}
\begin{canon}{正方}
\canongroup{定义}
\canonrow{AI 的迅猛发展}{2022 年以来，机器生成文字、图像、音乐和声音的能力突然变强，而且人人可用}{生成式 AI 爆发}{AI 只是工具}{常识；标识办法 2025}
\canonrow{人类创作者}{把创造文艺作品当成自己一部分的人，靠它吃饭的算，业余在写在画的也算}{把创作当志业的人}{人人都是创作者；会打字就是创作者}{汉典“创作”；《著作权法实施条例》第三条}
\canonrow{存在的意义}{作用与重要性：这个世界还有多需要、多珍视由人来创作；创作还能给创作者自己多少东西}{人来创作这件事的分量}{意义就是只有人能给的东西}{汉典“意义”}
\canongroup{判准}
\canonrow{判准}{和 AI 爆发之前比，由人来创作这件事整体变重还是变轻；需要和珍视两半都算，看多数创作者}{两半都算，看多数}{只看珍视，不看需要}{中立判准}
\canonrow{承认线}{一部分执行型的订单确实在减少，但题目问的是意义，不是订单数}{我方承认执行受冲击}{AI 没有替代任何人}{Hui 等 2024}
\canongroup{论点}
\canonrow{人味成了稀缺品}{机器内容无限以后，社会开始专门区分人和机器，并且在能选的时候选人}{人的来源成了稀缺品；机器越多，社会越要把人认出来}{人类作品更贵、更赚钱}{Bellaiche 2023；Pew 2025；Thaler 2025}
\canonrow{立意回到人}{AI 接走画得像、写得顺，留给人的是想画什么、选哪一张；门槛降下来，有话要说的人更多了}{执行交给机器，立意回到人}{AI 让所有创作者都更好；甲方也是创作者（只作防守）}{Zhou \& Lee 2024；北京互联网法院 2023}
\canonrow{全场追问}{如果人的来源不重要，为什么法院、行业和七成以上的普通人，都坚持要把人和机器分开？}{编剧、作家和法院为什么都在抢着把它写清楚}{}{Pew 2025；WGA 2023}
\canongroup{申论分工}
\canonrow{二辩申论}{拆反方立论主干，推进论点一；新层是行业规矩：2023 年编剧合同、2025 年人类创作认证}{回应反二，再拆立论}{把立论再念一遍}{WGA 2023；Authors Guild 2025}
\canonrow{三辩申论}{处理反方最强点，推进论点二；新层是获赞集中度下降和 Randy Travis}{先处理最强点，再推进}{和二辩申论讲同一件事}{Zhou \& Lee 2024；CBS 2024}
\canongroup{数据}
\canonrow{标签实验}{同一批 AI 画，贴上人类标签，喜欢、美感、深刻、价值四项评分全都更高}{贴人类标签评价更高}{人类作品本来就画得更好}{Bellaiche 等 2023}
\canonrow{皮尤调查}{七成六的美国成年人认为，能分辨作品出自人还是 AI 很重要（76\%，5,023 人）}{七成以上想分辨}{全世界的人都在乎}{Pew 2025}
\canonrow{益普索调查}{八国九千人里，八成认为纯 AI 音乐应清楚标注，超过一半为分不出来感到不安（80\% / 52\%）}{八成要求标注}{大家愿意为人类作品多付钱}{Deezer/Ipsos 2025}
\canonrow{获赞研究}{五万多名用户、四百万件作品：产量多了四分之一，获赞概率高了一半（25\% / 50\%），获赞集中度下降}{产量多四分之一，获赞高一半}{AI 让作品更有新意}{Zhou \& Lee 2024}
\canonrow{写作实验}{AI 的点子让原本最不会写故事的人进步最大}{新手进步最大}{AI 让所有人都写得更好且更多样}{Doshi \& Hauser 2024}
\canongroup{案例}
\canonrow{Thaler 案}{2025 年美国上诉法院判定版权法里的作者必须是人，2026 年最高法院拒绝再审}{法院只认人当作者}{美国法律禁止用 AI 创作}{Authors Alliance 2025；NCC 2026}
\canonrow{标识办法}{我国从 2025 年 9 月起要求 AI 生成的内容必须标识}{AI 内容要标识}{标识是为了尊重创作者}{网信办 2025}
\canonrow{AI 文生图案}{北京互联网法院 2023 年认定：设计提示词、设置参数，作个性化的选择和安排，体现独创性智力投入，著作权归人}{法律认的是人的选择}{AI 作品都有版权}{北京互联网法院 2023}
\canonrow{Randy Travis}{2013 年中风后几乎失语，2024 年团队用 AI 从 42 段旧录音还原嗓音，本人和老制作人把关，发了新歌}{借 AI 重新唱歌的人}{AI 替他唱了歌}{CBS 2024}
\canonrow{行业规矩}{2023 年好莱坞编剧新合同：AI 不能写、也不能改写剧本；2025 年美国作家协会推出人类创作认证}{编剧合同、人类创作认证}{所有作家都拿到了认证}{WGA 2023；Authors Guild 2025}
\end{canon}

\consubsection{10.2 反方口径表}
\begin{canon}{反方}
\canongroup{定义}
\canonrow{AI 的迅猛发展}{2022 年以来，机器生成文字、图像、音乐和声音的能力突然变强，而且可以不经过人直接交出成品}{生成式 AI 爆发}{AI 会取代所有创作者}{常识；标识办法 2025}
\canonrow{人类创作者}{把创造文艺作品当成自己一部分的人，职业的、业余的都算；尤其要看不出名、靠作品被需要的大多数}{不出名的大多数}{只有职业创作者才算}{汉典“创作”；《著作权法实施条例》第三条}
\canonrow{存在的意义}{作用与重要性：这个世界还有多需要、多珍视由人来创作；创作还能给创作者自己留下多少}{人来创作这件事的分量}{意义就是收入}{汉典“意义”}
\canongroup{判准}
\canonrow{判准}{和 AI 爆发之前比，由人来创作这件事整体变重还是变轻；需要和珍视两半都算，看多数创作者}{两半都算，但要看落在谁身上}{只看需要，不看珍视}{中立判准}
\canonrow{承认线}{法律和标注确实把人和机器分开了，门槛也确实降了；但被区分不等于被需要}{我方承认区分是真的}{法律和标注毫无意义}{Thaler 2025；标识办法 2025}
\canongroup{论点}
\canonrow{不再被需要}{对大多数不出名的创作者，被需要就是他存在意义里最实的一块；AI 让这份需要绕开了人，连最好的那批也挡不住}{需要在多数人身上缩水}{创作者都失业了；AI 已经取代了人}{Hui 等 2024；英国作家协会 2024}
\canonrow{手艺被掏空}{留下来的人，作品里属于他自己的那部分正在变少；构思去了甲方，执行给了机器，手里只剩修补}{作品越来越不像某一个人}{用 AI 的都不算创作者}{Zhou \& Lee 2024；Rest of World 2023}
\canonrow{夹击句}{用 AI，作品里的你变少；不用 AI，需要你的人变少}{两头都在变轻}{}{论点一 + 论点二}
\canonrow{全场追问}{正方说的被珍视，有多少落到了一个不出名的插画师身上？}{珍视落在谁身上}{}{Rosen 1981}
\canongroup{申论分工}
\canonrow{二辩申论}{拆正方立论主干，推进论点一；新层是顶尖自由职业者受冲击更大、配音字幕译者五成六的收入面临风险}{全场第一篇申论}{把立论再念一遍}{Hui 等 2024；CISAC 2024}
\canonrow{三辩申论}{回应正三，处理正方最强点，推进论点二；新层是重庆工作室十五裁五、广东画师“更高效也更累”}{先回应，再处理最强点}{和二辩申论讲同一件事}{Rest of World 2023}
\canongroup{数据}
\canonrow{平台研究}{ChatGPT 之后，写作类月收入降了百分之五以上，图像类降了将近一成；月订单降 2\%；顶尖者受冲击更大（5.2\% / 9.4\%）}{收入降百分之五以上}{自由职业者都没活了}{Hui 等 2024}
\canonrow{作家协会调查}{两成六的插画师、三成六的译者已经因为生成式 AI 丢了活（26\% / 36\%，787 份）}{四分之一的插画师丢了活}{英国一半创作者失业}{英国作家协会 2024}
\canonrow{CISAC 预测}{到 2028 年，音乐与视听创作者两成四、两成一的收入面临风险，配音字幕译者五成六（24\% / 21\% / 56\%）}{预测收入面临风险}{已经损失了五成六}{CISAC 2024}
\canonrow{新颖度}{五万多名用户、四百万件作品：平均内容新颖度下降，视觉新颖度从最好的到平均的全都在下降}{作品越来越像}{AI 让所有作品都变差}{Zhou \& Lee 2024}
\canonrow{盲听}{八国九千人的盲听里，九成七分不出人和 AI（97\%）}{九成七分不出}{没人在乎是不是人做的}{Deezer/Ipsos 2025}
\canonrow{AI 曲目}{一个音乐平台 2026 年 6 月高峰时，每天约九万首纯 AI 曲目，超过新上传的一半}{每天约九万首}{AI 歌曲占了一半播放量}{Deezer 2026}
\canongroup{案例}
\canonrow{Amber Yu}{曾花一个星期画一个穿传统服装舞狮的女子，一张海报收三千到七千元；后来只被请去修光线、修画歪的手脚，价钱是原来的十分之一}{修图，十分之一的价}{她失业了}{Rest of World 2023}
\canonrow{重庆与广东}{重庆一家工作室十五个角色设计师裁掉五个，“两个人能干原来十个人的活”；广东画师说 AI 让她们“更高效，也更累”}{十五裁五}{全中国画师都被裁了}{Rest of World 2023}
\canonrow{AI 声音案}{配音师的声音被 AI 化做成产品出售，北京互联网法院 2024 年判赔 25 万元}{声音被做成产品卖}{配音行业消失了}{中新网 2024}
\canonrow{猎头估计}{杭州一位游戏猎头说插画岗位一年降了约 70\%（单人估计，只进防守）}{有猎头估计约七成}{行业统计显示降了七成}{Rest of World 2023}
\end{canon}

\section{十一、论据出处清单}
\begin{sourcelist}
\source{1}{“意义”的词典义：作用与重要性}{汉典“意义”词条，2026 年访问，\url{https://www.zdic.net/hans/意义}}{作用与重要性}{2026-09-30}{已核实}{双方定义}{}
\source{2}{“创作”的词典义：出于己意的创造}{汉典“创作”词条（《国语辞典》），2026 年访问，\url{https://www.zdic.net/hans/创作}}{特指不事模仿、出于己意的创造。多指文学、艺术等作品的创造而言}{2026-09-30}{已核实}{双方定义}{}
\source{3}{《著作权法实施条例》第三条对“创作”的定义}{国务院，《中华人民共和国著作权法实施条例》（2013 修订），中国江苏网 2024 转载，\url{https://jsnews.jschina.com.cn/zt2024/fzsxxxpt/hdxf/xzfg/202404/t20240419_3393487.shtml}}{著作权法所称创作，是指直接产生文学、艺术和科学作品的智力活动}{2026-09-30}{已核实}{双方定义}{}
\source{4}{ChatGPT 后写作类自由职业者月订单 -2\%、月收入 -5.2\%；图像类 -3.7\% / -9.4\%}{Hui, Reshef \& Zhou (2024), Organization Science 35(6), DOI 10.1287/orsc.2023.18441；数字经 WashU Olin 新闻稿 2023-08 核实，\url{https://olin.washu.edu/about/news-and-media/news/2023/08/study-ai-tools-cause-a-decline-in-freelance-work-and-incomeat-least-in-the-short-run.php}}{monthly jobs for writing-related freelancers on Upwork declined by 2\%, while monthly earnings declined by 5.2\%}{2026-09-30}{已核实}{反方论点一、正方反驳库}{}
\source{5}{顶尖自由职业者受冲击更大}{同上，CESifo Working Paper No. 10601（2023）摘要，\url{https://ideas.repec.org/p/ces/ceswps/_10601.html}}{we find suggestive evidence that top freelancers are disproportionately affected by AI}{2026-09-30}{已核实}{反方二辩申论}{}
\source{6}{英国作家协会 2024 调查：26\% 插画师、36\% 译者已丢活；65\% 小说作者担心收入；787 份回应}{Society of Authors 2024 年 1 月调查，经欧洲作家理事会全文转载，\url{https://europeanwriterscouncil.eu/soa-survey-uk-ai-2024/}}{A quarter of illustrators (26\%) and over a third of translators (36\%) have already lost work due to generative AI}{2026-09-30}{已核实}{反方论点一}{}
\source{7}{CISAC 预测：到 2028 年音乐 24\%、视听 21\% 收入面临风险，配音字幕译者 56\%}{CISAC / PMP Strategy，2024-12 新闻稿，\url{https://www.cisac.org/Newsroom/news-releases/global-economic-study-shows-human-creators-future-risk-generative-ai}}{music and audiovisual creators will see respectively 24\% and 21\% of their revenues at risk of loss by 2028}{2026-09-30}{已核实}{反方论点一、反方二辩申论}{}
\source{8}{Deezer：2026 年 6 月高峰每天约 90,000 首 AI 曲目，超过新上传一半；只占播放 1--3\%；最多 85\% 播放为欺诈}{Deezer Newsroom，2026-07-21，\url{https://newsroom-deezer.com/2026/07/ai-music-exceeds-50-percent-daily-uploads-deezer/}}{90,000 AI-generated tracks per day now represent over 50\% of all new music uploads on Deezer}{2026-09-30}{已核实}{反方三辩申论、正方反驳库}{}
\source{9}{Deezer 不向用户推荐纯 AI 曲目（解释播放份额低的原因之一）}{Deezer Newsroom，2025-06-20，\url{https://newsroom-deezer.com/2025/06/deezer-launches-worlds-first-ai-tagging-system-for-music-streaming/}}{Deezer is currently excluding fully AI generated tracks from algorithmic and editorial recommendations}{2026-09-30}{已核实}{双方反驳库}{}
\source{10}{Deezer/益普索 8 国 9,000 人：97\% 盲听分不出；52\% 感到不安；80\% 要求标注}{Deezer Newsroom，2025-11-12，\url{https://newsroom-deezer.com/2025/11/deezer-ipsos-survey-ai-music/}}{80\% agree that 100\% AI-generated music should be clearly labeled}{2026-09-30}{已核实}{正方论点一；反方三辩申论}{}
\source{11}{非专业读者辨认 AI 诗准确率 46.6\%，更常把 AI 诗当成人写的}{Porter \& Machery (2024), Scientific Reports, DOI 10.1038/s41598-024-76900-1；摘要经约翰内斯堡大学机构库核实，\url{https://pure.uj.ac.za/en/publications/ai-generated-poetry-is-indistinguishable-from-human-written-poetr/}}{participants performed below chance levels in identifying AI-generated poems (46.6\% accuracy}{2026-09-30}{已核实}{反方反驳库、正方反驳库}{}
\source{12}{贴“人类创作”标签的同一批 AI 画，四项评分全高}{Bellaiche 等 (2023), Cognitive Research: Principles and Implications, DOI 10.1186/s41235-023-00499-6；Europe PMC 摘要，\url{https://europepmc.org/search?query=10.1186/s41235-023-00499-6}}{Study 1 found increased positive judgements for human- compared to AI-labelled art across all criteria}{2026-09-30}{已核实}{正方论点一}{}
\source{13}{5 万多名用户、400 万件作品：产量 +25\%，获赞概率 +50\%；平均内容新颖度下降，视觉新颖度全面下降；获赞集中度下降}{Zhou \& Lee (2024), PNAS Nexus, DOI 10.1093/pnasnexus/pgae052}{text-to-image AI significantly enhances human creative productivity by 25\% and increases the value}{2026-09-30}{已核实}{正方论点二；反方论点二}{}
\source{14}{AI 点子让最不会写的人进步最大，但故事更趋同}{Doshi \& Hauser (2024), Science Advances, DOI 10.1126/sciadv.adn5290}{generative AI-enabled stories are more similar to each other than stories by humans alone}{2026-09-30}{已核实}{正方论点二；反方论点二}{}
\source{15}{Thaler v. Perlmutter：作者必须是人（2025-03-18）；最高法院拒绝调卷（2026-03-02）}{Authors Alliance 2025-03-19，\url{https://www.authorsalliance.org/2025/03/19/thaler-v-perlmutter-d-c-court-of-appeals-confirms-that-a-non-human-machine-cannot-be-an-author-under-the-u-s-copyright-act/}；National Constitution Center 2026，\url{https://constitutioncenter.org/blog/supreme-court-denies-artificial-intelligence-authorship-claim-for-artwork-copyright}}{All of these statutory provisions collectively identify an 'author' as a human being}{2026-09-30}{已核实}{正方论点一}{}
\source{16}{北京互联网法院 AI 文生图案：提示词与参数的个性化选择安排体现独创性智力投入}{北京市人大常委会网站，2024-01-23，\url{https://www.bjrd.gov.cn/zyfb/zt/16j2crdh2024/bgjd/lygzbg/202401/t20240123_3543208.html}}{能够体现出原告的独创性智力投入，享有涉案图片的著作权}{2026-09-30}{已核实}{正方论点二}{}
\source{17}{全国首例 AI 生成声音人格权侵权案，判赔 25 万元}{中新网，2024-04-25，\url{https://www.chinanews.com.cn/sh/2024/04-25/10205621.shtml}}{赔偿原告各项损失25万元，并书面赔礼道歉}{2026-09-30}{已核实}{反方论点一}{}
\source{18}{Amber Yu：一张海报 3,000--7,000 元，后来只修图，价钱是原来的十分之一；舞狮插画画了一周；重庆工作室 15 裁 5；广东画师“更高效也更累”；猎头估计插画岗位一年降约 70\%}{Rest of World，Viola Zhou，2023-04-11，\url{https://restofworld.org/2023/ai-image-china-video-game-layoffs/}}{Amber Yu used to make 3,000 to 7,000 yuan (\$430 to \$1,000) for every video game poster she drew}{2026-09-30}{已核实}{反方论点二、反方三辩申论、结辩}{}
\source{19}{Randy Travis 2013 年中风，2024 年借 AI 从 42 段旧录音还原嗓音发新歌}{CBS News，2024-05-05，\url{https://www.cbsnews.com/news/randy-travis-sings-again-courtesy-of-ai-where-that-came-from/}}{In 2013 Travis suffered a massive stroke}{2026-09-30}{已核实}{正方三辩申论、正一结辩}{}
\source{20}{WGA 2023 合同：AI 不能写或改写剧本，公司不能要求编剧用 AI}{WGA 2023 MBA 官方摘要，\url{https://www.wgacontract2023.org/the-campaign/summary-of-the-2023-wga-mba}}{AI can't write or rewrite literary material}{2026-09-30}{已核实}{正方二辩申论}{}
\source{21}{《人工智能生成合成内容标识办法》2025 年 9 月 1 日起施行；答记者问写明为提醒用户辨别虚假信息}{网信办，2025-03-14，\url{https://www.cac.gov.cn/2025-03/14/c_1743654685899683.htm}；答记者问 \url{https://www.cac.gov.cn/2025-03/14/c_1743654685896173.htm}}{通过标识提醒用户辨别虚假信息}{2026-09-30}{已核实}{正方论点一；反方反驳库}{}
\source{22}{Authors Guild 人类创作认证（2025）：AI 生成的书越来越像人写的}{Authors Guild，2025-01-29，\url{https://authorsguild.org/news/ag-launches-human-authored-certification-to-preserve-authenticity-in-literature/}}{AI-generated books are flooding online marketplaces and increasingly look like, and sometimes even read like, human-authored books}{2026-09-30}{已核实}{正方二辩申论；反方反驳库}{}
\source{23}{皮尤 2025：76\% 认为能分辨 AI 或人很重要；53\% 认为 AI 会削弱人的创造性思考；5,023 人}{Pew Research Center，2025-09-17，\url{https://www.pewresearch.org/science/2025/09/17/how-americans-view-ai-and-its-impact-on-people-and-society/}}{76\% say it's extremely or very important to be able to tell if pictures, videos and text were made by AI or people}{2026-09-30}{已核实}{正方论点一}{}
\source{24}{Clarkesworld 2023 年 2 月收到 500 多份机器写的稿件，关闭投稿}{TechCrunch，2023-02-21，\url{https://techcrunch.com/2023/02/21/clarkesworld-ai-generated-submissions/}}{Clarkesworld had received more than 500 spam submissions in the month of February alone}{2026-09-30}{已核实}{双方反驳库备用}{}
\source{25}{原作价值来自“一次独特的创作行为”}{Newman \& Bloom (2012), JEP: General, DOI 10.1037/a0026035}{the assessment of the art object as a unique creative act (performance)}{2026-09-30}{已核实}{正方论点一学理}{}
\source{26}{努力启发式及其复制：复制结果混合，出价一项未复制}{Kruger 等 (2004), JESP 40: 91--98；Ziano 等 (2023), Collabra: Psychology, DOI 10.1525/collabra.87489}{We report two replications with mixed results}{2026-09-30}{已核实}{正方论点一边界}{}
\source{27}{信息丰富导致注意力贫乏}{Simon (1971), Designing Organizations for an Information-Rich World；Conversable Economist 2015 全文引录，\url{https://conversableeconomist.blogspot.com/2015/08/economics-of-information-overload.html}}{a wealth of information creates a poverty of attention}{2026-09-30}{已核实}{反方淘汰论点备用}{}
\source{28}{超级明星经济：少数人拿走巨额收入}{Rosen (1981), American Economic Review 71(5)，\url{https://pdodds.w3.uvm.edu/files/papers/others/1981/rosen1981a.pdf}}{relatively small numbers of people earn enormous amounts of money and dominate the activities in which they engage}{2026-09-30}{已核实}{反方论点一学理}{}
\source{29}{劳动过程理论：构思与执行分离、去技能化}{Braverman (1974), Labor and Monopoly Capital, Monthly Review Press}{}{}{有把握}{反方论点二学理}{查原书第五章（科学管理）关于 separation of conception from execution 的原句}
\source{30}{Porter \& Machery 实验二：告诉被试“人写的”评分更高}{azoai 2024-11-19 转述}{}{}{待核实}{备用（未进稿件）}{打开 Scientific Reports 原文 Study 2 结果段}
\source{31}{Authors Guild 人类创作认证 2026 年扩展到所有美国作者、非会员每本收费}{搜索摘要，未打开原页}{}{}{待核实}{备用（未进稿件）}{authorsguild.org/news/human-authored-certification-expands-to-all-authors}
\source{32}{Hui 等研究的观察窗口长度}{未核实}{}{}{待核实}{备用（未进稿件）}{Organization Science 原文数据部分}
\end{sourcelist}

\section{十二、备赛建议}
\begin{fields}
\begin{timeline}\when{第 1--2 天}{全队对齐定义、判准和两个论点标签；一辩背稿，二三辩背申论的固定段。}\when{第 3--4 天}{按记录卡练套件组装：每人至少把自己申论的每个模块念一遍，计时。}\when{第 5 天}{质询与被质询对练，双边计时，专门练“不打断也问得完”。}\when{第 6 天}{整场模辩，换边再打一场；赛后按口径表查滑坡。}\end{timeline}
\field{模辩重点检验}{\sub{正方}{被问“标签能带来订单吗”时是否守住“珍视不是价格”；是否主动承认执行型订单在减少。}\sub{反方}{是否滑成“AI 取代创作者”；是否把 CISAC 预测说成已经发生。}\sub{双方}{两篇申论有没有讲成同一件事；被质询时每答是否在八秒内。}}
\begin{checklist}\item 我方定义：人类创作者、存在的意义（两层）\item 判准：两半都算，看多数\item 两个论点标签与各自一条核心数据\item 申论分工：二辩拆立论加论点一，三辩处理最强点加论点二\item 全场追问那一句\end{checklist}
\end{fields}
\end{document}
